\bookchapter{Chapitre 28 : Dépenser juste pour acheter le bon — Guide ultime d'achat de smartphone}

\chapterintro{Pour acheter un smartphone, beaucoup de gens regardent des tests pendant une semaine et comparent les prix pendant un mois, pour finir par dépenser cinq ou six mille yuans dans un modèle phare qu'ils n'utilisent pas vraiment. Ce chapitre ne liste pas de tableaux de paramètres — vous les trouverez sur les sites officiels. Ce que je veux faire, c'est vous donner une « méthode de choix », afin qu'à l'avenir, quel que soit le nouveau smartphone sur le marché, vous puissiez juger par vous-même s'il vous convient.}

\booksection{Méthode de choix en trois étapes : besoins budget modèle}

Beaucoup de gens demandent en ligne « quel smartphone acheter pour 3000 yuans », une question à laquelle il est impossible de répondre — car on ne sait pas à quoi va principalement servir le smartphone. L'ordre de réflexion correct pour acheter un smartphone est :

\begin{tcbinfo}{Les trois étapes du choix}
\begin{enumerate}
  \item \textbf{D'abord, réfléchissez à ce qui compte le plus pour vous}: Photo ? Jeux ? Autonomie ? Écran ? Fluidité du système ? Ou seulement téléphoner et envoyer des messages WeChat ? Classez vos besoins par ordre de priorité.
  \item \textbf{Ensuite, fixez un plafond budgétaire}: Pas « environ combien », mais « au maximum combien je peux dépenser ». Fixez un plafond ferme, pour ne pas être tenté par des modèles de plus en plus chers.
  \item \textbf{Enfin, choisissez le modèle}: Avec votre classement de besoins et votre budget, trouvez le modèle correspondant.
\end{enumerate}

\end{tcbinfo}

\booksection{I. Comment lire les paramètres clés — inutile de tout comprendre, ces cinq suffisent}

Les tableaux de paramètres contiennent souvent des dizaines de données, mais l'utilisateur ordinaire n'a besoin de prêter attention qu'à ces cinq dimensions :

\begin{center}
\begin{tabularx}{\linewidth}{|X|X|X|}
\hline
\textbf{Dimension} & \textbf{Comment évaluer} & \textbf{En une phrase} \tabularnewline
\hline
Processeur (SoC) & Détermine pendant combien d'années le smartphone restera fluide. Norme 2026 : Snapdragon 8 Elite Gen5/Dimensity 9500 et au-dessus = haut de gamme ; Snapdragon 8s Gen4/Dimensity 8400 = milieu-haut de gamme ; Snapdragon 7+ Gen4/Dimensity 7350 = suffisant pour trois ans. En dessous de la série Snapdragon 6, prudence. & Achetez neuf, pas ancien ; le SoC détermine combien d'années votre smartphone restera fluide \tabularnewline
Écran & L'OLED offre de meilleures couleurs et consomme moins que le LCD ; un taux de rafraîchissement de 120 Hz est bien plus fluide que 60 Hz pour le défilement ; la résolution doit être au moins 1080p. La modulation PWM à haute fréquence (1920 Hz et au-dessus) est plus respectueuse des yeux. & OLED + 120 Hz + modulation PWM haute fréquence = le minimum pour un bon écran \tabularnewline
Caméra & \textbf{Ne regardez pas les mégapixels, mais la taille du capteur}. La taille du capteur compte bien plus que le nombre de pixels. L'ultra grand-angle est pratique ; le téléobjectif (zoom optique 3x et au-dessus) pour les sujets éloignés ; l'objectif macro est souvent gadget. & Grand capteur principal + téléobjectif pratique > plusieurs objectifs gadgets \tabularnewline
Batterie et charge & 4500 mAh est le minimum, 5000 mAh et au-dessus assurent une charge par jour sans problème. Pour la charge rapide, regardez les watts : 60 W et au-dessus c'est rapide, 100 W et au-dessus c'est très rapide. La charge sans fil est pratique mais lente. & 5000 mAh + charge rapide 60 W = combinaison sans anxiété \tabularnewline
Engagement de mises à jour système & Détermine pendant combien d'années le smartphone restera sécurisé. Apple offre généralement 5-6 ans de mises à jour système ; les phares Samsung 4 ans de système + 5 ans de sécurité ; Google Pixel 7 ans ; les modèles milieu de gamme nationaux ont des engagements plus courts (2-3 ans). & Pour utiliser longtemps, choisissez une marque avec un long engagement de mises à jour \tabularnewline
\hline
\end{tabularx}
\end{center}

\booksection{II. Recommandations par tranche budgétaire (basé sur le marché mi-2026)}

Les recommandations suivantes sont basées sur le marché chinois continental mi-2026, les prix sont les prix de lancement officiels. Les prix réels peuvent varier, référez-vous au prix au moment de l'achat.

\booksubsection{Moins de 1000 yuans : téléphoner, WeChat, ça suffit}

\textbf{Pour qui}: Étudiants, smartphone de secours, utilisateurs peu exigeants, principalement pour téléphoner et WeChat.

\textbf{Ce que vous pouvez acheter}: Redmi Note, Honor X, realme série numérique en configuration milieu-bas. Processeur Snapdragon série 6 ou Dimensity 6000, écran principalement LCD, photo correcte mais mise au point lente, fluidité système moyenne.

\textbf{Conseil d'achat}: À ce prix, regardez deux choses : la batterie est-elle assez grande (choisissez 5000 mAh ou plus) et la mémoire est-elle suffisante (choisissez au moins 6 Go + 128 Go, n'achetez pas 4 Go car ce sera lent). Redmi a l'avantage le plus net dans cette tranche.

\booksubsection{1000-2000 yuans : le meilleur rapport qualité-prix, la zone idéale pour 3-4 ans sans ralentissement}

\textbf{Pour qui}: Étudiants, budget limité mais envie de garder longtemps, pas de recherche extrême en photo et jeux.

\textbf{Modèles représentatifs}: Redmi Turbo, iQOO Z, Honor X haut de gamme, OnePlus Ace entrée de gamme.

\textbf{Conseil d'achat}: Choisissez un modèle équipé d'un Snapdragon 7+ Gen3 ou Dimensity 7300 et au-dessus, avec 8 Go + 256 Go. Cette configuration suffit largement pour un usage quotidien et occasionnellement Honor of Kings, et peut tenir 3-4 ans sans ralentissement.

\booksubsection{2000-3500 yuans : le cœur de gamme, performances et photo nettement améliorées}

\textbf{Pour qui}: Utilisateurs ordinaires souhaitant de bonnes photos, des jeux fluides et une longue durée de vie. C'est la tranche de prix la plus achetée.

\textbf{Modèles représentatifs}: Série numérique Xiaomi (comme le Xiaomi 17 standard aprèsréduction), vivo S haut de gamme, OPPO Reno, série numérique Honor, Redmi K.

\textbf{Conseil d'achat}: La caractéristique principale de cette tranche est le « smartphone polyvalent » — sans défaut majeur. Par rapport aux modèles à moins de 2000 yuans, les améliorations sont : appareil photo nettement meilleur (avec stabilisation optique), écran OLED de haute qualité, charge plus rapide (80 W et au-dessus courant), finition plus soignée. Choisissez 12 Go + 256 Go ou 12 Go + 512 Go, pour 4-5 ans sans problème.

\booksubsection{3500-5000 yuans : quasi-phare, suffisant pour tout le monde sauf les geeks}

\textbf{Pour qui}: Utilisateurs exigeants en photo, souhaitant une finition haut de gamme.

\textbf{Modèles représentatifs}: Série numérique Xiaomi standard, vivo X, OPPO Find X standard, anciens phares iPhone (comme l'iPhone 16).

\textbf{Conseil d'achat}: À ce prix, vous pouvez déjà acheter les phares de l'année précédente. Mon conseil :\textbf{Achetez le phare de l'année dernière, pas le milieu de gamme de cette année}. Les matériaux, la finition et le matériel photo du phare de l'année précédente sont généralement un cran au-dessus du milieu de gamme de cette année. À condition d'accepter qu'il ne soit pas « le dernier modèle ».

\booksubsection{Plus de 5000 yuans : zone phare, pour les budgets sans limite ou les exigences extrêmes}

\textbf{Pour qui}: Passionnés de photo, joueurs intensifs, utilisateurs cherchant un écran et une finition haut de gamme, créateurs de contenu vidéo sur smartphone.

\textbf{Modèles représentatifs}: iPhone 17 Pro/Pro Max, Samsung S26 Ultra, Huawei Pura 80 Ultra, Xiaomi 17 Ultra, vivo X300 Pro.

\textbf{Conseil d'achat}: À ce prix, les écarts entre marques sont de plus en plus faibles, le choix dépend surtout du système que vous préférez. Pour la photo, choisissez vivo ou Huawei ; pour la vidéo, l'iPhone ; pour l'écran et le S Pen, Samsung ; pour l'équilibre, Xiaomi. Honnêtement, à ce prix,\textbf{vous pouvez suivre votre cœur, choisissez celui qui vous plaît}.

\booksection{III. Choix par scénario — quels sont vos besoins ?}

Voici quelques « profils d'utilisation » courants, trouvez le vôtre :

\begin{center}
\begin{tabularx}{\linewidth}{|X|X|X|}
\hline
\textbf{Scénario d'utilisation} & \textbf{À quoi prêter attention} & \textbf{Direction recommandée} \tabularnewline
\hline
Passionné de photo & Grand capteur principal (1 pouce idéal), stabilisation optique, téléobjectif (zoom optique 3x et au-dessus), puce image & vivo X, Huawei Pura/Mate, Xiaomi Ultra \tabularnewline
Joueur & Processeur haut de gamme, taux d'échantillonnage tactile élevé, système de dissipation thermique, grande batterie + charge rapide & iQOO numérique, ROG Phone, Red Magic, Redmi K \tabularnewline
Autonomie prioritaire & Batterie de 5500 mAh et au-dessus, processeur économe, écran LTPO & Honor Magic, OnePlus, Redmi Note version grande batterie \tabularnewline
Seniors & Mode simplifié, grande police, son fort, grande batterie, nombreux points de SAV & Huawei Enjoy, Honor Play, OPPO A, Redmi Note \tabularnewline
Étudiants & Bon rapport qualité-prix, durable, grande batterie, contrôle parental & Redmi Note/Turbo, realme GT, iQOO Neo \tabularnewline
Utilisateurs écosystème Apple & Déjà équipé Mac/iPad/AirPods, l'iPhone est la seule option & iPhone 17 standard (rapport qualité-prix) ou Pro (expérience complète) \tabularnewline
Voyage à l'étranger & Prise en charge eSIM, couverture complète des bandes mondiales, services Google & iPhone, Samsung S, Pixel, OnePlus \tabularnewline
\hline
\end{tabularx}
\end{center}

\booksection{IV. Cinq pièges fréquents — vérifiez avant d'acheter}

\begin{enumerate}
  \item \textbf{Piège 1 : quatre ou cinq objectifs ne signifient pas bonne photo}— Les fabricants aiment empiler les objectifs dans les paramètres, dont deux sont souvent des « objectifs gadget » de 2 mégapixels (macro + profondeur), totalement inutiles. Pour évaluer la capacité photo d'un smartphone,\textbf{regardez seulement principal + ultra grand-angle + téléobjectif}ces trois suffisent. Les autres peuvent être ignorés.
  \item \textbf{Piège 2 : ne vous laissez pas intimider par 100 mégapixels}— Un nombre de pixels élevé ne signifie pas bonne qualité d'image. La taille du capteur, l'ouverture et le traitement algorithmique sont bien plus importants. Beaucoup de modèles phares utilisent un grand capteur de 50 mégapixels, et les photos sont bien meilleures que celles de certains modèles à 100 mégapixels à 1000 yuans.
  \item \textbf{Piège 3 : n'achetez pas de smartphone opérateur}— « Déposez du crédit, recevez un smartphone » semble avantageux, mais les smartphones opérateurs ont des mises à jour système lentes, des applications préinstallées impossibles à supprimer, et sont généralement verrouillés (seule la carte SIM de cet opérateur fonctionne). Apparemment économie, en réalité « verrouillage ».
  \item \textbf{Piège 4 : « l'écran courbe c'est haut de gamme » — c'est du marketing, pas une vérité}— L'écran courbe était autrefois le signe des modèles phares, mais de plus en plus de gens découvrent ses problèmes : protection d'écran difficile, fausses manips sur les bords, réparation coûteuse en cas de casse. La tendance 2026 est de revenir à « légèrement courbé » voire « totalement plat ». Ne choisissez pas un écran courbe sous prétexte que « courbe = haut de gamme ».
  \item \textbf{Piège 5 : tous les « smartphones IA » ne sont pas vraiment utiles}— En 2026, tous les fabricants crient « smartphone IA », mais beaucoup se contentent de renommer leur assistant vocal existant. Les fonctions IA vraiment utiles sont : la traduction en temps réel, l'effacement IA (supprimer des personnes d'une photo), la retouche IA, l'inférence de grand modèle côté appareil, le résumé intelligent local. Les autres fonctions voyantes sont surtout des gadgets.
\end{enumerate}

\booksection{V. Quand acheter au meilleur prix ?}

\begin{itemize}
  \item \textbf{618 (fin mai à mi-juin)}: La période de l'année avec les plus grandes baisses de prix. De nombreux fabricants baissent les prix de 800-1500 yuans pour booster les ventes. C'est la meilleure fenêtre d'achat du premier semestre.
  \item \textbf{Double 11 (fin octobre à mi-novembre)}: La grande promotion du second semestre, d'une ampleur similaire au 618. À cette période, de nombreuses nouvelles générations de phares viennent de sortir, et la baisse de prix de la génération précédente est maximale.
  \item \textbf{Un mois après le lancement}: Au lancement, prix catalogue ; un mois plus tard, le prix des canaux (boutiques tierces) est généralement 200-500 yuans moins cher.
  \item \textbf{Deux semaines avant le Nouvel An chinois à la Fête des Lanternes}: Le « festival des achats du Nouvel An » des magasins physiques pour écouler les stocks, il est possible de négocier avec le vendeur.
  \item \textbf{Moments déconseillés}: Acheter le jour même de la présentation — à moins de vouloir vraiment être parmi les premiers, le prix de lancement est le plus élevé.
\end{itemize}

\booksection{VI. Quel canal choisir — où acheter en confiance ?}

\begin{itemize}
  \item \textbf{Boutique officielle de la marque / boutique phare Tmall / JD direct}: Le plus sûr, pas d'avantage de prix mais pas de souci de reconditionné ou de contrefaçon, meilleur SAV.
  \item \textbf{« Subvention de 10 milliards » de Pinduoduo}: Le prix le plus bas, mais choisissez des boutiques avec le label « Subvention de 10 milliards ». N'achetez pas chez les boutiques sans label. Filmez le déballage, en cas de problème vous pourrez retourner.
  \item \textbf{Magasins physiques}: Vous pouvez toucher et essayer. Attention : le vendeur poussera les modèles à forte marge (pas les plus adaptés), il est conseillé de cibler 2-3 modèles en ligne avant d'entrer, et de dire « je ne veux voir que ceux-ci ».
  \item \textbf{Plateformes d'occasion (Xianyu/Zhuanzhuan)}: Moins cher mais risqué. Pour acheter en occasion, retenez trois règles : demandez la facture ou la capture de commande, demandez le numéro IMEI pour vérification,\textbf{ne traitez que via la plateforme, ne virez jamais en privé}.
\end{itemize}

\begin{tcbpractice}{Exercice : choisissez un smartphone avec cette méthode}
\begin{enumerate}
  \item Prenez une feuille, notez vos deux besoins clés (par exemple : autonomie d'abord, photo ensuite).
  \item Fixez un plafond budgétaire (par exemple 3500 yuans).
  \item Ouvrez JD ou Tmall, filtrez par prix, sélectionnez 3-5 modèles.
  \item Cherchez sur Bilibili « modèle XX utilisation d'un mois » (pas le test du premier jour, mais le ressenti après au moins un mois), regardez deux ou trois vidéos.
  \item Confirmez la fiabilité du canal, commandez, et filmez le déballage.
\end{enumerate}

\end{tcbpractice}

\begin{tcbsummary}{Résumé du chapitre}
\begin{itemize}
  \item Trois étapes du choix : besoins budget modèle
  \item Cinq paramètres clés : processeur, écran, caméra, batterie/charge, engagement de mises à jour
  \item 1500-3500 yuans est la zone idéale avec le meilleur rapport qualité-prix
  \item Acheter le phare de l'année précédente est plus avantageux que le milieu de gamme de cette année
  \item Cinq pièges : arnaque des multicapteurs, piège des mégapixels, smartphone opérateur, mythe de l'écran courbe, gadget IA
  \item Le 618 et le Double 11 sont les deux meilleures fenêtres d'achat de l'année
\end{itemize}

\end{tcbsummary}
